% ============================================================
% Chapter 2: Related Work — OUTLINE (Phase 1)
% 各節末に必ず「本研究との相違」を置く（Kotoba Tech ルール）。
% ============================================================
\section{Component Technologies for Automated Physical Model Building}
\begin{itemize}
  \item \textbf{P1.} Automating physical model building requires a chain of component technologies, several of which have been developed in our group.
  \begin{itemize}
    \item Document collection, formula extraction, and variable-equivalence judgment (ProcessBERT\cite{ProcessBERT}) precede model assembly.
    % cite: 数式抽出・文書情報抽出の先行研究（Phase 2で実文献を検証して追加。捏造禁止）
    \item (difference) These works produce the equation database; this study addresses the downstream step of selecting a solvable equation set from it.
  \end{itemize}
\end{itemize}

\section{Information Retrieval and Reranking}
\begin{itemize}
  \item \textbf{P2.} Classical information retrieval ranks items independently by query--item relevance, and modern systems refine the top candidates with a second-stage reranker.
  \begin{itemize}
    % cite: TF-IDF/BM25 と retrieve-then-rerank の代表文献（Phase 2で検証して追加）
    \item Lexical methods (TF-IDF) and learned rerankers assume that item relevances are independent given the query.
    \item Diversification methods consider inter-item redundancy but not solvability.
    \item (difference) Our items are interdependent: a correct answer is a set that closes a system of equations, which motivates set-aware features.
  \end{itemize}
\end{itemize}

\section{Direct Generation of Domain Knowledge by Large Language Models}
\begin{itemize}
  \item \textbf{P3.} Large language models can generate domain equations directly from a description, which makes them a natural alternative to retrieval.
  \begin{itemize}
    % cite: LLMの科学知識生成・hallucination に関する代表文献（Phase 2で検証して追加）
    \item Generated equations may be physically valid yet differ in notation, and empirical equations with fitted coefficients are hard to reproduce verbatim.
    \item (difference) We use an LLM only as an external baseline (and for generating case descriptions); our method retrieves from a database, which keeps provenance and reproducibility.
  \end{itemize}
\end{itemize}

\section{Mathematical Information Retrieval}
\begin{itemize}
  \item \textbf{P4.} Mathematical information retrieval searches for single formulas similar to a query formula or text.
  \begin{itemize}
    % cite: math IR / formula search の代表文献（Phase 2で検証して追加）
    \item Existing systems rank individual formulas by structural or textual similarity.
    \item (difference) Our task returns a set of equations that jointly form a solvable model, a requirement absent from formula search.
  \end{itemize}
\end{itemize}
